\documentclass[10pt,a4paper]{article}

% ----- Encodage & langue -----
\usepackage[T1]{fontenc}
\usepackage[utf8]{inputenc}
\usepackage[french]{babel}

% ----- Mise en page -----
\usepackage{geometry}
\geometry{top=1.5cm,bottom=1.5cm,left=1.5cm,right=1.5cm}

% Colonnes
\usepackage{multicol}
\setlength{\columnsep}{0.75cm}            % espace entre colonnes
\setlength{\columnseprule}{0.2pt}      % épaisseur de la ligne
\def\columnseprulecolor{\color{Couleur8}} % couleur de la ligne

% Maths et figures
\usepackage{amsmath, amssymb}
\usepackage{tikz}
\usetikzlibrary{babel,arrows,calc,fit,patterns,plotmarks,shapes.geometric,shapes.misc,shapes.symbols,shapes.arrows,shapes.callouts, shapes.multipart, shapes.gates.logic.US,shapes.gates.logic.IEC, er, automata,backgrounds,chains,topaths,trees,petri,mindmap,matrix, calendar, folding,fadings,through,positioning,scopes,decorations.fractals,decorations.shapes,decorations.text,decorations.pathmorphing,decorations.pathreplacing,decorations.footprints,decorations.markings,shadows}

\usetikzlibrary{calc}

% Listes compactes
\usepackage{enumitem}
\setlist{nosep, itemsep=0.4em, parsep=0pt}
\setlist[enumerate,1]{label=\alph*.}

% Couleurs
\usepackage{xcolor}
\definecolor{exoblue}{HTML}{005F73}

%----------------------------------------------------------------------------------------
%	 COULEURS
%----------------------------------------------------------------------------------------

\definecolor{Couleur1}{HTML}{8b1f30}
\definecolor{Couleur2}{HTML}{725E74}
\definecolor{Couleur3}{HTML}{78985b}
\definecolor{Couleur4}{HTML}{db3e56}
\definecolor{Couleur5}{HTML}{487C4B}
\definecolor{Couleur6}{HTML}{CF6876}
\definecolor{Couleur7}{HTML}{A5C03F} %Vert
\definecolor{Couleur8}{HTML}{D36740} %Orange
\definecolor{Couleur9}{HTML}{7BB48A} %Bleu-Vert
\definecolor{Couleur10}{HTML}{EA8556} %Orange
\definecolor{Couleur11}{HTML}{E6B459} %Jaune
\definecolor{Couleur12}{HTML}{9AC8EB} %Bleu

\definecolor{Bord1}{HTML}{302635}
\definecolor{Bord2}{HTML}{725E74}
\definecolor{Bord3}{HTML}{938999}
\definecolor{Bord4}{HTML}{817888}
\definecolor{Bord5}{HTML}{487C4B}
\definecolor{Bord6}{HTML}{8C436B}
\definecolor{Bord7}{HTML}{A5C03F} %Vert
\definecolor{Bord8}{HTML}{D36740} %Orange

% ----- Police sans serif -----
\renewcommand{\familydefault}{\sfdefault}

% ----- Exercice -----
\newcounter{exo}
\newcommand{\exo}[1][]{%
  \refstepcounter{exo}%
  \textbf{\textcolor{Couleur8}{\\[0.3cm]\theexo.}}%
  \if\relax\detokenize{#1}\relax\else\ \textbf{#1}\fi\ %
}

\begin{document}
\begin{Large}\textbf{\textcolor{Couleur8}{EXERCICES - Chapitre transversal - Résolution de problèmes}} \end{Large}

%https://coopmaths.fr/alea/?uuid=7fb24&id=6N5-3&n=10&d=10&s=1&s2=11&s3=true&cd=1&alea=q20i&uuid=4e89b&id=6N4A-1&n=1&d=10&s=2&s2=4&s3=3&cd=1&alea=V7Ih&uuid=4e89b&id=6N4A-1&n=1&d=10&s=2&s2=4&s3=2&cd=1&alea=wlKP&uuid=4e89c&id=6N4A-2&n=4&d=10&s=1&s2=5&cd=1&alea=w6gw


% @see : https://coopmaths.fr/alea?uuid=7fb24&id=6N5-3&n=10&d=10&s=1&s2=11&s3=true&alea=q20i&cd=1&cols=1
\exo Dans chaque problème, coche les informations qui \textbf{servent} à sa résolution.

\begin{enumerate}[itemsep=2em]
	\item Un cargo mesurant 109 m transporte 74 gros conteneurs de 26 tonnes chacun de Marseille à Rio de Janeiro. Ce bateau transporte aussi 13 petits conteneurs pour une masse totale de 182 tonnes.\\Quelle est la masse de chacun des petits conteneurs, sachant qu'ils ont tous la même masse ? 
    
    	$\square\;$ 109 m\qquad $\square\;$ 74 conteneurs\qquad $\square\;$ 26 tonnes\qquad $\square\;$ 13 conteneurs\qquad $\square\;$ 182 tonnes\qquad 
	\item Le village de Sainte-Farida-Les-Trois-Vallées compte 872 habitants et se situe à une altitude de 280 m. À 7 km de là, le village de Saint-Mehdi-Le-Bouquetin, situé 134 m plus haut, compte 337 habitants de moins.\\Combien de garçons compte le village de Sainte-Farida-Les-Trois-Vallées ? 
    
    	$\square\;$ 337 habitants\qquad $\square\;$ 7 km\qquad $\square\;$ 134 m\qquad $\square\;$ 872 habitants\qquad $\square\;$ 280 m\qquad 
	\item Sébastien, un élève de 6ème, de 11 ans, mesure 1,34 m. Fiona a 5 ans de plus que Sébastien et mesure 31 cm de plus.\\En quelle classe est Fiona ? 
    
    	$\square\;$ 11 ans\qquad $\square\;$ 1,34 m\qquad $\square\;$ 31 cm\qquad $\square\;$ 5 ans\qquad $\square\;$ 6ème\qquad 
	\item Le frère de Julien, âgé de 45 ans, se rend 4 fois par semaine à Toulouse en train. Une fois arrivé, il prend le métro à 7 h 45, après avoir acheté systématiquement le même journal, dans un kiosque de la gare, qui coûte 1,40 €. Son trajet en métro dure 45 minutes pour se rendre au travail.\\À quelle heure le frère de Julien arrive-t-il à son travail ? 
    
    	$\square\;$ 1,40 €\qquad $\square\;$ 45 ans\qquad $\square\;$ 45 min\qquad $\square\;$ 7 h 45\qquad $\square\;$ 4 fois\qquad 
	\item La grand-mère de Océane lui a acheté un superbe vélo de 15 vitesses, coûtant 576 €, avec des roues de 18 pouces. Pour la protéger, son père lui a offert un casque et du matériel d'éclairage valant 28,85 €. La grand-mère de Océane a décidé de payer le vélo en 2 fois.\\Quel est le montant de chaque versement que payera la grand-mère de Océane ? 
    
    	$\square\;$ 28,85 €\qquad $\square\;$ 15 vitesses\qquad $\square\;$ 576 €\qquad $\square\;$ 18 pouces\qquad $\square\;$ 2 fois\qquad 
	\item Vanessa décide de programmer la box de sa tante pour enregistrer un film prévu le dimanche 24 janvier et une émission prévue le lendemain. Le film doit commencer à 16 h 25 et se terminer à 18 h 40. L'émission commence à 17 h 50 et dure 42 minutes.\\À quelle heure de sa journée, Vanessa décide-t-elle de programmer la box de sa tante ? 
    
    	$\square\;$ dimanche 24 janvier\qquad $\square\;$ 42 minutes\qquad $\square\;$ 17 h 50\qquad $\square\;$ 18 h 40\qquad $\square\;$ 16 h 25\qquad 
	\item Au marché, Zacharie achète 5 barquettes de haricots verts de 750\,g chacune à 3,84\, €~pièce  et 4\,ananas coûtant 2,67\, €~l'unité.\\Quel est le prix total des boissons achetées ? 
    
    	$\square\;$ 750 g\qquad $\square\;$ 3,84 €\qquad $\square\;$ 5 barquettes\qquad $\square\;$ 2,67 €\qquad $\square\;$ 4 ananas\qquad 
	\item Un livreur part de son entrepôt avec 27 colis. Au premier arrêt, le plus près, il dépose 18 colis. 13 km plus loin, il livre le reste de ses colis. Ensuite, à 12 h 35, le livreur reprend la même route et retourne à l'entrepôt, à 28 km de là.\\Combien de colis le livreur a-t-il déposé à son deuxième arrêt ? 
    
    	$\square\;$ 12 h 35\qquad $\square\;$ 27 colis\qquad $\square\;$ 18 colis\qquad $\square\;$ 13 km\qquad $\square\;$ 28 km\qquad 
	\item Dans une classe de 24 élèves âgés de 14  à 16  ans, un professeur distribue à chaque enfant 10 livres pesant 440 g chacun.\\Quel est, dans cette classe, le nombre exact d'enfants de 15 ans ? 
    
    	$\square\;$ 440 g\qquad $\square\;$ 14 ans\qquad $\square\;$ 16 ans\qquad $\square\;$ 24 élèves\qquad $\square\;$ 10 livres\qquad 
	\item Tania vient de lire en 1 h 30 un manga qu'elle avait payé 6,64 €. Elle a remarqué que sur chaque page, il y avait exactement 8 cases. C'est grâce au billet de 10 €~que lui a donné son père, que Tania a pu s'acheter ce livre de 84 pages.\\À quelle heure Tania a-t-elle commencé à lire son manga ? 
    
    	$\square\;$ 1 h 30\qquad $\square\;$ 84 pages\qquad $\square\;$ 6,64 €\qquad $\square\;$ 8 cases\qquad $\square\;$ 10 €\qquad 
\end{enumerate}
 
 % Colonnes

\setlength{\columnsep}{0.75cm}            % espace entre colonnes
\setlength{\columnseprule}{0pt}      % épaisseur de la ligne
\def\columnseprulecolor{\color{Couleur8}} % couleur de la ligne

% @see : https://coopmaths.fr/alea?uuid=4e89b&id=6N4A-1&n=1&d=10&s=2&s2=4&s3=3&alea=V7Ih&cd=1&cols=1
\exo Associer chaque problème avec le schéma correspondant.
\begin{multicols}{4}
\noindent \begin{tikzpicture}[baseline,scale = 0.3]

    \tikzset{
      point/.style={
        thick,
        draw,
        cross out,
        inner sep=0pt,
        minimum width=5pt,
        minimum height=5pt,
      },
    }
    \clip (-2,-1) rectangle (12.5,6);
    	\draw [color={Couleur10},line width = 1] (0,0)--(12,0)--(12,4)--(0,4)--cycle;
	\draw [color={Couleur10}] (0,2)--(12,2);
	\draw [color={Couleur10}] (6,2)--(6,4);
	\draw [color={black}] (6,1) node[anchor = center,scale=1, rotate = 0] {?};
	\draw [color={black}] (3,3) node[anchor = center,scale=1, rotate = 0] {40};
	\draw [color={black}] (9,3) node[anchor = center,scale=1, rotate = 0] {6};
	\draw [color={black}] (-1,4) node[anchor = center,scale=1, rotate = 0] {A.};

\end{tikzpicture}\\
\begin{tikzpicture}[baseline,scale = 0.3]

    \tikzset{
      point/.style={
        thick,
        draw,
        cross out,
        inner sep=0pt,
        minimum width=5pt,
        minimum height=5pt,
      },
    }
    \clip (-2,-1) rectangle (12.5,6);
    	\draw [color={Couleur10},line width = 1] (0,0)--(12,0)--(12,4)--(0,4)--cycle;
	\draw [color={Couleur10}] (0,2)--(12,2);
	\draw [color={Couleur10}] (2,2)--(2,4);
	\draw [color={black}] (6,1) node[anchor = center,scale=1, rotate = 0] {240};
	\draw [color={black}] (1,3) node[anchor = center,scale=1, rotate = 0] {40};
	\draw [color={black}] (3,3) node[anchor = center,scale=1, rotate = 0] {40};
	\draw[color ={black},<->] (0,4.5)--(12,4.5);
	\draw [color={Couleur10}] (4,2)--(4,4);
	\draw [color={Couleur10}] (10,2)--(10,4);
	\draw [color={black}] (7,3) node[anchor = center,scale=1, rotate = 0] {...};
	\draw [color={black}] (11,3) node[anchor = center,scale=1, rotate = 0] {40};
	\draw [color={black}] (6,5) node[anchor = center,scale=1, rotate = 0] {?};
	\draw [color={black}] (-1,4) node[anchor = center,scale=1, rotate = 0] {E.};

\end{tikzpicture}\\
\begin{tikzpicture}[baseline,scale = 0.3]

    \tikzset{
      point/.style={
        thick,
        draw,
        cross out,
        inner sep=0pt,
        minimum width=5pt,
        minimum height=5pt,
      },
    }
    \clip (-2,-1) rectangle (12.5,6);
    	\draw [color={Couleur10},line width = 1] (0,0)--(12,0)--(12,4)--(0,4)--cycle;
	\draw [color={Couleur10}] (0,2)--(12,2);
	\draw [color={Couleur10}] (2,2)--(2,4);
	\draw [color={black}] (6,1) node[anchor = center,scale=1, rotate = 0] {48};
	\draw [color={black}] (1,3) node[anchor = center,scale=1, rotate = 0] {6};
	\draw [color={black}] (3,3) node[anchor = center,scale=1, rotate = 0] {6};
	\draw[color ={black},<->] (0,4.7)--(12,4.7);
	\draw [color={Couleur10}] (4,2)--(4,4);
	\draw [color={Couleur10}] (10,2)--(10,4);
	\draw [color={black}] (7,3) node[anchor = center,scale=1, rotate = 0] {. . .};
	\draw [color={black}] (11,3) node[anchor = center,scale=1, rotate = 0] {6};
	\draw [color={black}] (6,5.2) node[anchor = center,scale=1, rotate = 0] {?};
	\draw [color={black}] (-1,4) node[anchor = center,scale=1, rotate = 0] {B.};

\end{tikzpicture}\\
\begin{tikzpicture}[baseline,scale = 0.3]

    \tikzset{
      point/.style={
        thick,
        draw,
        cross out,
        inner sep=0pt,
        minimum width=5pt,
        minimum height=5pt,
      },
    }
    \clip (-2,-1) rectangle (12.5,6);
    	\draw [color={Couleur10},line width = 1] (0,0)--(12,0)--(12,4)--(0,4)--cycle;
	\draw [color={Couleur10}] (0,2)--(12,2);
	\draw [color={Couleur10}] (6,2)--(6,4);
	\draw [color={black}] (6,1) node[anchor = center,scale=1, rotate = 0] {40};
	\draw [color={black}] (3,3) node[anchor = center,scale=1, rotate = 0] {?};
	\draw [color={black}] (9,3) node[anchor = center,scale=1, rotate = 0] {6};
	\draw [color={black}] (-1,4) node[anchor = center,scale=1, rotate = 0] {F.};

\end{tikzpicture}\\
\begin{tikzpicture}[baseline,scale = 0.3]

    \tikzset{
      point/.style={
        thick,
        draw,
        cross out,
        inner sep=0pt,
        minimum width=5pt,
        minimum height=5pt,
      },
    }
    \clip (-2,-1) rectangle (12.5,6);
    	\draw [color={Couleur10},line width = 1] (0,0)--(12,0)--(12,4)--(0,4)--cycle;
	\draw [color={Couleur10}] (0,2)--(12,2);
	\draw [color={Couleur10}] (6,2)--(6,4);
	\draw [color={black}] (6,1) node[anchor = center,scale=1, rotate = 0] {?};
	\draw [color={black}] (3,3) node[anchor = center,scale=1, rotate = 0] {48};
	\draw [color={black}] (9,3) node[anchor = center,scale=1, rotate = 0] {40};
	\draw [color={black}] (-1,4) node[anchor = center,scale=1, rotate = 0] {C.};

\end{tikzpicture}\\
\begin{tikzpicture}[baseline,scale = 0.3]

    \tikzset{
      point/.style={
        thick,
        draw,
        cross out,
        inner sep=0pt,
        minimum width=5pt,
        minimum height=5pt,
      },
    }
    \clip (-2,-1) rectangle (12.5,6);
    	\draw [color={Couleur10},line width = 1] (0,0)--(12,0)--(12,4)--(0,4)--cycle;
	\draw [color={Couleur10}] (0,2)--(12,2);
	\draw [color={Couleur10}] (6,2)--(6,4);
	\draw [color={black}] (6,1) node[anchor = center,scale=1, rotate = 0] {48};
	\draw [color={black}] (3,3) node[anchor = center,scale=1, rotate = 0] {40};
	\draw [color={black}] (9,3) node[anchor = center,scale=1, rotate = 0] {?};
	\draw [color={black}] (-1,4) node[anchor = center,scale=1, rotate = 0] {G.};

\end{tikzpicture}\\
\begin{tikzpicture}[baseline,scale = 0.3]

    \tikzset{
      point/.style={
        thick,
        draw,
        cross out,
        inner sep=0pt,
        minimum width=5pt,
        minimum height=5pt,
      },
    }
    \clip (-2,-1) rectangle (12.5,6);
    	\draw [color={Couleur10},line width = 1] (0,0)--(12,0)--(12,4)--(0,4)--cycle;
	\draw [color={Couleur10}] (0,2)--(12,2);
	\draw [color={Couleur10}] (2,2)--(2,4);
	\draw [color={black}] (6,1) node[anchor = center,scale=1, rotate = 0] {?};
	\draw [color={black}] (1,3) node[anchor = center,scale=1, rotate = 0] {40};
	\draw [color={black}] (3,3) node[anchor = center,scale=1, rotate = 0] {40};
	\draw[color ={black},<->] (0,4.5)--(12,4.5);
	\draw [color={Couleur10}] (4,2)--(4,4);
	\draw [color={Couleur10}] (10,2)--(10,4);
	\draw [color={black}] (7,3) node[anchor = center,scale=1, rotate = 0] {. . .};
	\draw [color={black}] (11,3) node[anchor = center,scale=1, rotate = 0] {40};
	\draw [color={black}] (6,5) node[anchor = center,scale=1, rotate = 0] {6};
	\draw [color={black}] (-1,4) node[anchor = center,scale=1, rotate = 0] {D.};

\end{tikzpicture}\\
\begin{tikzpicture}[baseline,scale = 0.3]

    \tikzset{
      point/.style={
        thick,
        draw,
        cross out,
        inner sep=0pt,
        minimum width=5pt,
        minimum height=5pt,
      },
    }
    \clip (-2,-1) rectangle (12.5,6);
    	\draw [color={Couleur10},line width = 1] (0,0)--(12,0)--(12,4)--(0,4)--cycle;
	\draw [color={Couleur10}] (0,2)--(12,2);
	\draw [color={Couleur10}] (2,2)--(2,4);
	\draw [color={black}] (6,1) node[anchor = center,scale=1, rotate = 0] {48};
	\draw [color={black}] (1,3) node[anchor = center,scale=1, rotate = 0] {?};
	\draw [color={black}] (3,3) node[anchor = center,scale=1, rotate = 0] {?};
	\draw[color ={black},<->] (0,4.7)--(12,4.7);
	\draw [color={Couleur10}] (4,2)--(4,4);
	\draw [color={Couleur10}] (10,2)--(10,4);
	\draw [color={black}] (7,3) node[anchor = center,scale=1, rotate = 0] {. . .};
	\draw [color={black}] (11,3) node[anchor = center,scale=1, rotate = 0] {?};
	\draw [color={black}] (6,5.2) node[anchor = center,scale=1, rotate = 0] {6};
	\draw [color={black}] (-1,4) node[anchor = center,scale=1, rotate = 0] {H.};

\end{tikzpicture}


    \end{multicols}

% Colonnes

\setlength{\columnsep}{0.75cm}            % espace entre colonnes
\setlength{\columnseprule}{0.2pt}      % épaisseur de la ligne
\def\columnseprulecolor{\color{Couleur8}} % couleur de la ligne


    \begin{multicols}{2}
\begin{enumerate}[itemsep=1em]
	\item Kamel a besoin de $240$ gommes.\\Il en récupère $40$ chaque jour.\\Au bout de combien de temps aura-t-il le nécessaire ?
	\item J'ai payé $48$ €~pour des gâteaux coûtant $6$ €~chacun.\\Combien en ai-je acheté ?
	\item $6$ photos identiques coûtent $48$ €.\\Quel est le prix d'une d'entre elles ?
	\item Dans un sac, il y a $40$ cartes de v\oe ux et dans l'autre, il y en a $48$.\\Combien y en a-t-il de plus dans ce deuxième sac ?
	\item Elsa avait $40$ livres lundi. \\Le lendemain, elle en a trouvé $6$ autres.\\Combien cela lui en fait-il ?
	\item Yvette a $40$ ans.\\Sachant qu'elle a $6$ ans de plus que son frère, quel âge a celui-ci ?
	\item Roxane récupère $48$ billes dans une salle, puis $40$ dans une autre.\\Combien en a-t-elle en tout ?
	\item Hélène récupère $6$ paquets de $40$ cahiers chacun.\\Combien en a-t-elle en tout ?
\end{enumerate}
 
    \end{multicols}

% @see : https://coopmaths.fr/alea?uuid=4e89b&id=6N4A-1&n=1&d=10&s=2&s2=4&s3=2&alea=wlKP&cd=1&cols=1
\exo Déterminer un schéma à associer à chaque problème. Il n'est pas demandé de résoudre le problème.
\begin{multicols}{2}
\begin{enumerate}[itemsep=1em]
	\item François a $7$ ans de moins que sa s\oe ur Julie.\\Sachant qu'il a $36$ ans, quel âge a sa s\oe ur ?
	\item Karole a acheté $7$ stylos à $36$ €~pièce.\\Combien a-t-elle payé ?
	\item Guillaume a besoin de $252$ gommes.\\Il en récupère $36$ chaque jour.\\Au bout de combien de temps aura-t-il le nécessaire ?
	\item Yvette a acheté $36$ cahiers pour les donner à ses amis.\\Il lui en reste encore $7$ à donner.\\Combien en a-t-elle déjà distribué ?
	\columnbreak
	\item J'ai $63$ cahiers dans mon sac et je dois les regrouper par $7$.\\Combien puis-je faire de tas ?
	\item J'ai $63$ photos dans mon sac et je souhaite les partager avec mes $6$ amis.\\Quelle sera la part de chacun ?
	\item Dans un sac, il y a $36$ cartes de v\oe ux et dans l'autre, il y en a $63$.\\Combien y en a-t-il de plus dans ce deuxième sac ?
	\item Un lot de crayons coûte $63$ €~et un lot de cartes de v\oe ux coûte $36$ €.\\Combien coûte l'ensemble ?
\end{enumerate}
 \end{multicols}

\newpage
% @see : https://coopmaths.fr/alea?uuid=4e89c&id=6N4A-2&n=4&d=10&s=1&s2=5&alea=w6gw&cd=1&cols=1
\exo Trouver les nombres manquants dans ces problèmes afin que énoncé, réponse et schéma se correspondent.
\begin{multicols}{2}
\begin{enumerate}[itemsep=1em]
	\item Un tee-shirt coûte $10$ €~et une paire de baskets coûte $109$ €. Béatrice a $60$ €. Combien Béatrice doit-elle avoir en plus pour acheter ces deux objets ?

\medskip

  Réponse :\\
  - Calcul du prix total en €~: $10+109=119$.\\
  - Calcul de la différence en €~: $119-60=59$.\\
  Béatrice doit avoir $59$ €~de plus pour acheter ces deux objets.

\medskip

  \begin{tikzpicture}[scale=0.8]
\draw[fill=Couleur10!15, draw=Couleur10, very thick] (0.0,2.2) rectangle (2.4,3.0);
\draw[fill=Couleur10!15, draw=Couleur10, very thick] (2.4,2.2) rectangle (8.0,3.0);
\draw[fill=Couleur10!15, draw=Couleur10, very thick] (0.0,0.9) rectangle (4.8,1.7);
\draw[<->,thick, draw=black] (4.8,1.5) -- (8.0,1.5);
\node[] at (6.4,1) {\textcolor{black}{\ldots}};
\draw[dashed, thick, draw=black] (8.0,0.9000000000000001) -- (8.0,1.7);
\end{tikzpicture}

\medskip

	\item Yann participe à un triathlon. Il nage $200$ m, fait ensuite du vélo et ensuite court sur $1\,000$ m. Il a parcouru en tout $4\,200$ m. Quelle distance a-t-il parcourue à vélo ?

\medskip

      Réponse :\\
      - Calcul de la distance parcourue en nage et en course à pied en m : $200+1\,000=1\,200$\\
      - Distance parcourue à vélo en m : $4\,200-1\,200=$ $3\,000$.

\medskip

  \begin{tikzpicture}[scale=0.8]
\draw[decorate,decoration={brace,amplitude=10pt},xshift=0pt,yshift=0pt,draw=black] (0.0,3) -- node[above=10pt, pos=0.5] {\textcolor{black}{\ldots}} (7.2,3);
\draw[fill=Couleur10!15, draw=Couleur10, very thick] (0.0,2.2) rectangle (2.4,3.0); 
\draw[fill=Couleur10!15, draw=Couleur10, very thick] (2.4,2.2) rectangle (4.8,3.0); 
\draw[fill=Couleur10!15, draw=Couleur10, very thick] (4.8,2.2) rectangle (7.2,3.0); 
\end{tikzpicture}

\medskip
\columnbreak
	\item Tania effectue à vélo une étape de montagne du tour de France. Elle grimpe trois cols qui ont les dénivelés suivants : $600$ m pour le premier col, $1\,200$ m pour le deuxième et $950$ m pour le dernier. Quelle dénivelé cumulé a-t-elle grimpé ?

\medskip

      Réponse : Tania a grimpé au total en m : $600+1\,200+950=$ $2\,750$.

\medskip

  \begin{tikzpicture}[scale=0.8]
\draw[decorate,decoration={brace,amplitude=10pt},xshift=0pt,yshift=0pt,draw=black] (0.0,3) -- node[above=10pt, pos=0.5] {\textcolor{black}{\ldots}} (7.2,3);
\draw[fill=Couleur10!15, draw=Couleur10, very thick] (0.0,2.2) rectangle (2.4,3.0) ;
\draw[fill=Couleur10!15, draw=Couleur10, very thick] (2.4,2.2) rectangle (4.8,3.0) ;
\draw[fill=Couleur10!15, draw=Couleur10, very thick] (4.8,2.2) rectangle (7.2,3.0) ;
\end{tikzpicture}

\medskip

	\item Karine a $5$ boîtes de chocolats chacun contenant $30$ chocolats. Quand elle les distribue à ses amis, chacun en a $37$ et il en reste 2. Combien a-t-elle d'amis ?

\medskip

    Réponse :\\
    - Calcul du nombre de chocolats total : $5 \times 30 = 150$.\\
    - Calcul du nombre de chocolats partagés : $150-2=148$.\\
    - Nombre d'amis : $148 \div 37 = 4$.\\
    Karine a $4$ amis.

\medskip

  \begin{tikzpicture}[scale=0.8]
\draw[decorate,decoration={brace,amplitude=10pt},xshift=0pt,yshift=0pt,draw=black] (0.0,3) -- node[above=10pt, pos=0.5] {\textcolor{black}{\ldots}} (8.0,3);
\draw[fill=Couleur10!15, draw=Couleur10, very thick] (0.0,2.2) rectangle (2.4,3.0);

\draw[fill=Couleur10!15, draw=Couleur10, very thick] (2.4,2.2) rectangle (8.0,3.0);

\draw[fill=Couleur10!15, draw=Couleur10, very thick] (0.0,0.9) rectangle (2.4,1.7);

\draw[fill=Couleur10!15, draw=Couleur10, very thick] (2.4,0.9) rectangle (5.6,1.7);

\draw[fill=Couleur10!15, draw=Couleur10, very thick] (5.6,0.9) rectangle (8.0,1.7);
\node[anchor=center, align=center] at (6.8,1.3) {\textcolor{black}{$2$}};
\draw[<->,thick, draw=black] (0.0,0.7000000000000002) -- (5.6,0.7000000000000002); \node[] at (2.8,0) {\textcolor{black}{\ldots}};
\end{tikzpicture}

\medskip

\end{enumerate}
 \end{multicols}

\newpage
\setcounter{exo}{0}
\begin{Large}\textbf{\textcolor{Couleur8}{Correction - Chapitre transversal - Résolution de problèmes}} \end{Large}

\exo
\vspace{-0.85cm}
\begin{multicols}{2}
\begin{enumerate}[itemsep=0.5em]
\item {\bfseries \color{Couleur10}13 conteneurs} et {\bfseries \color{Couleur10}182 tonnes} sont utiles pour la résolution du problème.\\La solution du problème est donnée par : {\bfseries \color{Couleur10}182 tonnes}$\,\div\,${\bfseries \color{Couleur10}13}.
\item Aucune donnée n'est utile pour la résolution du problème.\\On ne peut pas répondre à ce problème. Il manque des informations.
\item Aucune donnée n'est utile pour la résolution du problème.\\On ne peut pas répondre à ce problème. Il manque des informations.
\item {\bfseries \color{Couleur10}7 h 45} et {\bfseries \color{Couleur10}45 min} sont utiles pour la résolution du problème.\\La solution du problème est donnée par : {\bfseries \color{Couleur10}7 h 45}$\,+\,${\bfseries \color{Couleur10}45 min}.
\item {\bfseries \color{Couleur10}576 €} et {\bfseries \color{Couleur10}2 fois} sont utiles pour la résolution du problème.\\La solution du problème est donnée par : {\bfseries \color{Couleur10}576 €}$\,\div\,${\bfseries \color{Couleur10}2}.
\item Aucune donnée n'est utile pour la résolution du problème.\\On ne peut pas répondre à ce problème. Il manque des informations.
\item Aucune donnée n'est utile pour la résolution du problème.\\On ne peut pas répondre à ce problème. Il manque des informations.
\item {\bfseries \color{Couleur10}27 colis} et {\bfseries \color{Couleur10}18 colis} sont utiles pour la résolution du problème.\\La solution du problème est donnée par : {\bfseries \color{Couleur10}27 colis}$\,-\,${\bfseries \color{Couleur10}18 colis}.
\item Aucune donnée n'est utile pour la résolution du problème.\\On ne peut pas répondre à ce problème. Il manque des informations.
\item Aucune donnée n'est utile pour la résolution du problème.\\On ne peut pas répondre à ce problème. Il manque des informations.
\end{enumerate}
\end{multicols}
 

\exo
\vspace{-0.85cm}
\begin{multicols}{2}
 Kamel a besoin de $240$ gommes.\\Il en récupère $40$ chaque jour.\\Au bout de combien de temps aura-t-il le nécessaire ?

\medskip
Cet énoncé est associé avec le schéma {\bfseries \color{Couleur10}E}.\\\begin{tikzpicture}[baseline,scale = 0.3]

    \tikzset{
      point/.style={
        thick,
        draw,
        cross out,
        inner sep=0pt,
        minimum width=5pt,
        minimum height=5pt,
      },
    }
    \clip (-2,-1) rectangle (12.5,6);
    	\draw [color={Couleur10},line width = 1] (0,0)--(12,0)--(12,4)--(0,4)--cycle;
	\draw [color={Couleur10}] (0,2)--(12,2);
	\draw [color={Couleur10}] (2,2)--(2,4);
	\draw [color={black}] (6,1) node[anchor = center,scale=1, rotate = 0] {240};
	\draw [color={black}] (1,3) node[anchor = center,scale=1, rotate = 0] {40};
	\draw [color={black}] (3,3) node[anchor = center,scale=1, rotate = 0] {40};
	\draw[color ={black},<->] (0,4.5)--(12,4.5);
	\draw [color={Couleur10}] (4,2)--(4,4);
	\draw [color={Couleur10}] (10,2)--(10,4);
	\draw [color={black}] (7,3) node[anchor = center,scale=1, rotate = 0] {...};
	\draw [color={black}] (11,3) node[anchor = center,scale=1, rotate = 0] {40};
	\draw [color={black}] (6,5) node[anchor = center,scale=1, rotate = 0] {?};
	\draw [color={black}] (-1,4) node[anchor = center,scale=1, rotate = 0] {E.};

\end{tikzpicture}
 J'ai payé $48$ €~pour des gâteaux coûtant $6$ €~chacun.\\Combien en ai-je acheté ?

\medskip
Cet énoncé est associé avec le schéma {\bfseries \color{Couleur10}B}.\\\begin{tikzpicture}[baseline,scale = 0.3]

    \tikzset{
      point/.style={
        thick,
        draw,
        cross out,
        inner sep=0pt,
        minimum width=5pt,
        minimum height=5pt,
      },
    }
    \clip (-2,-1) rectangle (12.5,6);
    	\draw [color={Couleur10},line width = 1] (0,0)--(12,0)--(12,4)--(0,4)--cycle;
	\draw [color={Couleur10}] (0,2)--(12,2);
	\draw [color={Couleur10}] (2,2)--(2,4);
	\draw [color={black}] (6,1) node[anchor = center,scale=1, rotate = 0] {48};
	\draw [color={black}] (1,3) node[anchor = center,scale=1, rotate = 0] {6};
	\draw [color={black}] (3,3) node[anchor = center,scale=1, rotate = 0] {6};
	\draw[color ={black},<->] (0,4.7)--(12,4.7);
	\draw [color={Couleur10}] (4,2)--(4,4);
	\draw [color={Couleur10}] (10,2)--(10,4);
	\draw [color={black}] (7,3) node[anchor = center,scale=1, rotate = 0] {. . .};
	\draw [color={black}] (11,3) node[anchor = center,scale=1, rotate = 0] {6};
	\draw [color={black}] (6,5.2) node[anchor = center,scale=1, rotate = 0] {?};
	\draw [color={black}] (-1,4) node[anchor = center,scale=1, rotate = 0] {B.};

\end{tikzpicture}
 $6$ photos identiques coûtent $48$ €.\\Quel est le prix d'une d'entre elles ?

\medskip
Puisque le partage se fait avec mes $5$ amis, le partage se fait donc entre $6$ personnes (mes $5$ amis et moi). 

\medskip
Cet énoncé est associé avec le schéma {\bfseries \color{Couleur10}H}.\\\begin{tikzpicture}[baseline,scale = 0.3]

    \tikzset{
      point/.style={
        thick,
        draw,
        cross out,
        inner sep=0pt,
        minimum width=5pt,
        minimum height=5pt,
      },
    }
    \clip (-2,-1) rectangle (12.5,6);
    	\draw [color={Couleur10},line width = 1] (0,0)--(12,0)--(12,4)--(0,4)--cycle;
	\draw [color={Couleur10}] (0,2)--(12,2);
	\draw [color={Couleur10}] (2,2)--(2,4);
	\draw [color={black}] (6,1) node[anchor = center,scale=1, rotate = 0] {48};
	\draw [color={black}] (1,3) node[anchor = center,scale=1, rotate = 0] {?};
	\draw [color={black}] (3,3) node[anchor = center,scale=1, rotate = 0] {?};
	\draw[color ={black},<->] (0,4.7)--(12,4.7);
	\draw [color={Couleur10}] (4,2)--(4,4);
	\draw [color={Couleur10}] (10,2)--(10,4);
	\draw [color={black}] (7,3) node[anchor = center,scale=1, rotate = 0] {. . .};
	\draw [color={black}] (11,3) node[anchor = center,scale=1, rotate = 0] {?};
	\draw [color={black}] (6,5.2) node[anchor = center,scale=1, rotate = 0] {6};
	\draw [color={black}] (-1,4) node[anchor = center,scale=1, rotate = 0] {H.};

\end{tikzpicture}
 Dans un sac, il y a $40$ cartes de v\oe ux et dans l'autre, il y en a $48$.\\Combien y en a-t-il de plus dans ce deuxième sac ?

\medskip
Cet énoncé est associé avec le schéma {\bfseries \color{Couleur10}G}.\\\begin{tikzpicture}[baseline,scale = 0.3]

    \tikzset{
      point/.style={
        thick,
        draw,
        cross out,
        inner sep=0pt,
        minimum width=5pt,
        minimum height=5pt,
      },
    }
    \clip (-2,-1) rectangle (12.5,6);
    	\draw [color={Couleur10},line width = 1] (0,0)--(12,0)--(12,4)--(0,4)--cycle;
	\draw [color={Couleur10}] (0,2)--(12,2);
	\draw [color={Couleur10}] (6,2)--(6,4);
	\draw [color={black}] (6,1) node[anchor = center,scale=1, rotate = 0] {48};
	\draw [color={black}] (3,3) node[anchor = center,scale=1, rotate = 0] {40};
	\draw [color={black}] (9,3) node[anchor = center,scale=1, rotate = 0] {?};
	\draw [color={black}] (-1,4) node[anchor = center,scale=1, rotate = 0] {G.};

\end{tikzpicture}

\columnbreak

 Elsa avait $40$ livres lundi. \\Le lendemain, elle en a trouvé $6$ autres.\\Combien cela lui en fait-il ?


\medskip
Cet énoncé est associé avec le schéma {\bfseries \color{Couleur10}A}.\\\begin{tikzpicture}[baseline,scale = 0.3]

    \tikzset{
      point/.style={
        thick,
        draw,
        cross out,
        inner sep=0pt,
        minimum width=5pt,
        minimum height=5pt,
      },
    }
    \clip (-2,-1) rectangle (12.5,6);
    	\draw [color={Couleur10},line width = 1] (0,0)--(12,0)--(12,4)--(0,4)--cycle;
	\draw [color={Couleur10}] (0,2)--(12,2);
	\draw [color={Couleur10}] (6,2)--(6,4);
	\draw [color={black}] (6,1) node[anchor = center,scale=1, rotate = 0] {?};
	\draw [color={black}] (3,3) node[anchor = center,scale=1, rotate = 0] {40};
	\draw [color={black}] (9,3) node[anchor = center,scale=1, rotate = 0] {6};
	\draw [color={black}] (-1,4) node[anchor = center,scale=1, rotate = 0] {A.};

\end{tikzpicture}
 Yvette a $40$ ans.\\Sachant qu'elle a $6$ ans de plus que son frère, quel âge a celui-ci ?

\medskip
Cet énoncé est associé avec le schéma {\bfseries \color{Couleur10}F}.\\\begin{tikzpicture}[baseline,scale = 0.3]

    \tikzset{
      point/.style={
        thick,
        draw,
        cross out,
        inner sep=0pt,
        minimum width=5pt,
        minimum height=5pt,
      },
    }
    \clip (-2,-1) rectangle (12.5,6);
    	\draw [color={Couleur10},line width = 1] (0,0)--(12,0)--(12,4)--(0,4)--cycle;
	\draw [color={Couleur10}] (0,2)--(12,2);
	\draw [color={Couleur10}] (6,2)--(6,4);
	\draw [color={black}] (6,1) node[anchor = center,scale=1, rotate = 0] {40};
	\draw [color={black}] (3,3) node[anchor = center,scale=1, rotate = 0] {?};
	\draw [color={black}] (9,3) node[anchor = center,scale=1, rotate = 0] {6};
	\draw [color={black}] (-1,4) node[anchor = center,scale=1, rotate = 0] {F.};

\end{tikzpicture}
 Roxane récupère $48$ billes dans une salle, puis $40$ dans une autre.\\Combien en a-t-elle en tout ?

\medskip
Cet énoncé est associé avec le schéma {\bfseries \color{Couleur10}C}.\\\begin{tikzpicture}[baseline,scale = 0.3]

    \tikzset{
      point/.style={
        thick,
        draw,
        cross out,
        inner sep=0pt,
        minimum width=5pt,
        minimum height=5pt,
      },
    }
    \clip (-2,-1) rectangle (12.5,6);
    	\draw [color={Couleur10},line width = 1] (0,0)--(12,0)--(12,4)--(0,4)--cycle;
	\draw [color={Couleur10}] (0,2)--(12,2);
	\draw [color={Couleur10}] (6,2)--(6,4);
	\draw [color={black}] (6,1) node[anchor = center,scale=1, rotate = 0] {?};
	\draw [color={black}] (3,3) node[anchor = center,scale=1, rotate = 0] {48};
	\draw [color={black}] (9,3) node[anchor = center,scale=1, rotate = 0] {40};
	\draw [color={black}] (-1,4) node[anchor = center,scale=1, rotate = 0] {C.};

\end{tikzpicture}
 Hélène récupère $6$ paquets de $40$ cahiers chacun.\\Combien en a-t-elle en tout ?

\medskip
Cet énoncé est associé avec le schéma {\bfseries \color{Couleur10}D}.\\\begin{tikzpicture}[baseline,scale = 0.3]

    \tikzset{
      point/.style={
        thick,
        draw,
        cross out,
        inner sep=0pt,
        minimum width=5pt,
        minimum height=5pt,
      },
    }
    \clip (-2,-1) rectangle (12.5,6);
    	\draw [color={Couleur10},line width = 1] (0,0)--(12,0)--(12,4)--(0,4)--cycle;
	\draw [color={Couleur10}] (0,2)--(12,2);
	\draw [color={Couleur10}] (2,2)--(2,4);
	\draw [color={black}] (6,1) node[anchor = center,scale=1, rotate = 0] {?};
	\draw [color={black}] (1,3) node[anchor = center,scale=1, rotate = 0] {40};
	\draw [color={black}] (3,3) node[anchor = center,scale=1, rotate = 0] {40};
	\draw[color ={black},<->] (0,4.5)--(12,4.5);
	\draw [color={Couleur10}] (4,2)--(4,4);
	\draw [color={Couleur10}] (10,2)--(10,4);
	\draw [color={black}] (7,3) node[anchor = center,scale=1, rotate = 0] {. . .};
	\draw [color={black}] (11,3) node[anchor = center,scale=1, rotate = 0] {40};
	\draw [color={black}] (6,5) node[anchor = center,scale=1, rotate = 0] {6};
	\draw [color={black}] (-1,4) node[anchor = center,scale=1, rotate = 0] {D.};

\end{tikzpicture}
 \end{multicols}

\exo \vspace{-0.85cm}
\begin{multicols}{2}
 François a $7$ ans de moins que sa s\oe ur Julie.\\Sachant qu'il a $36$ ans, quel âge a sa s\oe ur ?

\medskip
Cet énoncé peut être associé avec le schéma ci-dessous.\\\begin{tikzpicture}[baseline,scale = 0.3]

    \tikzset{
      point/.style={
        thick,
        draw,
        cross out,
        inner sep=0pt,
        minimum width=5pt,
        minimum height=5pt,
      },
    }
    \clip (-1.5,-1) rectangle (12.5,6);
    	\draw [color={Couleur10},line width = 1] (0,0)--(12,0)--(12,4)--(0,4)--cycle;
	\draw [color={Couleur10}] (0,2)--(12,2);
	\draw [color={Couleur10}] (6,2)--(6,4);
	\draw [color={black}] (6,1) node[anchor = center,scale=1, rotate = 0] {?};
	\draw [color={black}] (3,3) node[anchor = center,scale=1, rotate = 0] {36};
	\draw [color={black}] (9,3) node[anchor = center,scale=1, rotate = 0] {7};
	\draw [color={black}] (-1,4) node[anchor = center,scale=1, rotate = 0] {};

\end{tikzpicture}
 Karole a acheté $7$ stylos à $36$ €~pièce.\\Combien a-t-elle payé ?

\medskip
Cet énoncé peut être associé avec le schéma ci-dessous.\\\begin{tikzpicture}[baseline,scale = 0.3]

    \tikzset{
      point/.style={
        thick,
        draw,
        cross out,
        inner sep=0pt,
        minimum width=5pt,
        minimum height=5pt,
      },
    }
    \clip (-1.5,-1) rectangle (12.5,6);
    	\draw [color={Couleur10},line width = 1] (0,0)--(12,0)--(12,4)--(0,4)--cycle;
	\draw [color={Couleur10}] (0,2)--(12,2);
	\draw [color={Couleur10}] (2,2)--(2,4);
	\draw [color={black}] (6,1) node[anchor = center,scale=1, rotate = 0] {?};
	\draw [color={black}] (1,3) node[anchor = center,scale=1, rotate = 0] {36};
	\draw [color={black}] (3,3) node[anchor = center,scale=1, rotate = 0] {36};
	\draw[color ={black},<->] (0,4.5)--(12,4.5);
	\draw [color={Couleur10}] (4,2)--(4,4);
	\draw [color={Couleur10}] (10,2)--(10,4);
	\draw [color={black}] (7,3) node[anchor = center,scale=1, rotate = 0] {. . .};
	\draw [color={black}] (11,3) node[anchor = center,scale=1, rotate = 0] {36};
	\draw [color={black}] (6,5) node[anchor = center,scale=1, rotate = 0] {7};
	\draw [color={black}] (-1,4) node[anchor = center,scale=1, rotate = 0] {};

\end{tikzpicture}
 Guillaume a besoin de $252$ gommes.\\Il en récupère $36$ chaque jour.\\Au bout de combien de temps aura-t-il le nécessaire ?

\medskip
Cet énoncé peut être associé avec le schéma ci-dessous.\\\begin{tikzpicture}[baseline,scale = 0.3]

    \tikzset{
      point/.style={
        thick,
        draw,
        cross out,
        inner sep=0pt,
        minimum width=5pt,
        minimum height=5pt,
      },
    }
    \clip (-1.5,-1) rectangle (12.5,6);
    	\draw [color={Couleur10},line width = 1] (0,0)--(12,0)--(12,4)--(0,4)--cycle;
	\draw [color={Couleur10}] (0,2)--(12,2);
	\draw [color={Couleur10}] (2,2)--(2,4);
	\draw [color={black}] (6,1) node[anchor = center,scale=1, rotate = 0] {252};
	\draw [color={black}] (1,3) node[anchor = center,scale=1, rotate = 0] {36};
	\draw [color={black}] (3,3) node[anchor = center,scale=1, rotate = 0] {36};
	\draw[color ={black},<->] (0,4.5)--(12,4.5);
	\draw [color={Couleur10}] (4,2)--(4,4);
	\draw [color={Couleur10}] (10,2)--(10,4);
	\draw [color={black}] (7,3) node[anchor = center,scale=1, rotate = 0] {...};
	\draw [color={black}] (11,3) node[anchor = center,scale=1, rotate = 0] {36};
	\draw [color={black}] (6,5) node[anchor = center,scale=1, rotate = 0] {?};
	\draw [color={black}] (-1,4) node[anchor = center,scale=1, rotate = 0] {};

\end{tikzpicture}
 Yvette a acheté $36$ cahiers pour les donner à ses amis.\\Il lui en reste encore $7$ à donner.\\Combien en a-t-elle déjà distribué ?

\medskip
Cet énoncé peut être associé avec le schéma ci-dessous.\\\begin{tikzpicture}[baseline,scale = 0.3]

    \tikzset{
      point/.style={
        thick,
        draw,
        cross out,
        inner sep=0pt,
        minimum width=5pt,
        minimum height=5pt,
      },
    }
    \clip (-1.5,-1) rectangle (12.5,6);
    	\draw [color={Couleur10},line width = 1] (0,0)--(12,0)--(12,4)--(0,4)--cycle;
	\draw [color={Couleur10}] (0,2)--(12,2);
	\draw [color={Couleur10}] (6,2)--(6,4);
	\draw [color={black}] (6,1) node[anchor = center,scale=1, rotate = 0] {36};
	\draw [color={black}] (3,3) node[anchor = center,scale=1, rotate = 0] {?};
	\draw [color={black}] (9,3) node[anchor = center,scale=1, rotate = 0] {7};
	\draw [color={black}] (-1,4) node[anchor = center,scale=1, rotate = 0] {};

\end{tikzpicture}

\columnbreak

 J'ai $63$ cahiers dans mon sac et je dois les regrouper par $7$.\\Combien puis-je faire de tas ?

\medskip
Cet énoncé peut être associé avec le schéma ci-dessous.\\\begin{tikzpicture}[baseline,scale = 0.3]

    \tikzset{
      point/.style={
        thick,
        draw,
        cross out,
        inner sep=0pt,
        minimum width=5pt,
        minimum height=5pt,
      },
    }
    \clip (-1.5,-1) rectangle (12.5,6);
    	\draw [color={Couleur10},line width = 1] (0,0)--(12,0)--(12,4)--(0,4)--cycle;
	\draw [color={Couleur10}] (0,2)--(12,2);
	\draw [color={Couleur10}] (2,2)--(2,4);
	\draw [color={black}] (6,1) node[anchor = center,scale=1, rotate = 0] {63};
	\draw [color={black}] (1,3) node[anchor = center,scale=1, rotate = 0] {7};
	\draw [color={black}] (3,3) node[anchor = center,scale=1, rotate = 0] {7};
	\draw[color ={black},<->] (0,4.7)--(12,4.7);
	\draw [color={Couleur10}] (4,2)--(4,4);
	\draw [color={Couleur10}] (10,2)--(10,4);
	\draw [color={black}] (7,3) node[anchor = center,scale=1, rotate = 0] {. . .};
	\draw [color={black}] (11,3) node[anchor = center,scale=1, rotate = 0] {7};
	\draw [color={black}] (6,5.2) node[anchor = center,scale=1, rotate = 0] {?};
	\draw [color={black}] (-1,4) node[anchor = center,scale=1, rotate = 0] {};

\end{tikzpicture}
 J'ai $63$ photos dans mon sac et je souhaite les partager avec mes $6$ amis.\\Quelle sera la part de chacun ?

\medskip
Puisque le partage se fait avec mes $6$ amis, le partage se fait donc entre $7$ personnes (mes $6$ amis et moi). 

\medskip
Cet énoncé peut être associé avec le schéma ci-dessous.\\\begin{tikzpicture}[baseline,scale = 0.3]

    \tikzset{
      point/.style={
        thick,
        draw,
        cross out,
        inner sep=0pt,
        minimum width=5pt,
        minimum height=5pt,
      },
    }
    \clip (-1.5,-1) rectangle (12.5,6);
    	\draw [color={Couleur10},line width = 1] (0,0)--(12,0)--(12,4)--(0,4)--cycle;
	\draw [color={Couleur10}] (0,2)--(12,2);
	\draw [color={Couleur10}] (2,2)--(2,4);
	\draw [color={black}] (6,1) node[anchor = center,scale=1, rotate = 0] {63};
	\draw [color={black}] (1,3) node[anchor = center,scale=1, rotate = 0] {?};
	\draw [color={black}] (3,3) node[anchor = center,scale=1, rotate = 0] {?};
	\draw[color ={black},<->] (0,4.7)--(12,4.7);
	\draw [color={Couleur10}] (4,2)--(4,4);
	\draw [color={Couleur10}] (10,2)--(10,4);
	\draw [color={black}] (7,3) node[anchor = center,scale=1, rotate = 0] {. . .};
	\draw [color={black}] (11,3) node[anchor = center,scale=1, rotate = 0] {?};
	\draw [color={black}] (6,5.2) node[anchor = center,scale=1, rotate = 0] {7};
	\draw [color={black}] (-1,4) node[anchor = center,scale=1, rotate = 0] {};

\end{tikzpicture}
 Dans un sac, il y a $36$ cartes de v\oe ux et dans l'autre, il y en a $63$.\\Combien y en a-t-il de plus dans ce deuxième sac ?

\medskip
Cet énoncé peut être associé avec le schéma ci-dessous.\\\begin{tikzpicture}[baseline,scale = 0.3]

    \tikzset{
      point/.style={
        thick,
        draw,
        cross out,
        inner sep=0pt,
        minimum width=5pt,
        minimum height=5pt,
      },
    }
    \clip (-1.5,-1) rectangle (12.5,6);
    	\draw [color={Couleur10},line width = 1] (0,0)--(12,0)--(12,4)--(0,4)--cycle;
	\draw [color={Couleur10}] (0,2)--(12,2);
	\draw [color={Couleur10}] (6,2)--(6,4);
	\draw [color={black}] (6,1) node[anchor = center,scale=1, rotate = 0] {63};
	\draw [color={black}] (3,3) node[anchor = center,scale=1, rotate = 0] {36};
	\draw [color={black}] (9,3) node[anchor = center,scale=1, rotate = 0] {?};
	\draw [color={black}] (-1,4) node[anchor = center,scale=1, rotate = 0] {};

\end{tikzpicture}
 Un lot de crayons coûte $63$ €~et un lot de cartes de v\oe ux coûte $36$ €.\\Combien coûte l'ensemble ?

\medskip
Cet énoncé peut être associé avec le schéma ci-dessous.\\\begin{tikzpicture}[baseline,scale = 0.3]

    \tikzset{
      point/.style={
        thick,
        draw,
        cross out,
        inner sep=0pt,
        minimum width=5pt,
        minimum height=5pt,
      },
    }
    \clip (-1.5,-1) rectangle (12.5,6);
    	\draw [color={Couleur10},line width = 1] (0,0)--(12,0)--(12,4)--(0,4)--cycle;
	\draw [color={Couleur10}] (0,2)--(12,2);
	\draw [color={Couleur10}] (6,2)--(6,4);
	\draw [color={black}] (6,1) node[anchor = center,scale=1, rotate = 0] {?};
	\draw [color={black}] (3,3) node[anchor = center,scale=1, rotate = 0] {63};
	\draw [color={black}] (9,3) node[anchor = center,scale=1, rotate = 0] {36};
	\draw [color={black}] (-1,4) node[anchor = center,scale=1, rotate = 0] {};

\end{tikzpicture}
 \end{multicols}

\newpage

\exo 
\vspace{-0.85cm}
\begin{multicols}{2}
\begin{enumerate}[itemsep=1em]
\item Le problème :

\medskip
Un tee-shirt coûte $10$ €~et une paire de baskets coûte $109$ €. Béatrice a $60$ €. Combien Béatrice doit-elle avoir en plus pour acheter ces deux objets ?

\medskip

  Réponse :\\
  - Calcul du prix total en €~: $10+109=119$.\\
  - Calcul de la différence en €~: $119-60=59$.\\
  Béatrice doit avoir $59$ €~de plus pour acheter ces deux objets.

\medskip

  se modélise par le schéma suivant :

\medskip

\begin{tikzpicture}[scale=0.8]
\draw[fill=Couleur10!15, draw=Couleur10, very thick] (0.0,2.2) rectangle (2.4,3.0) node[pos=0.5] {\textcolor{black}{${\color{Couleur10}\boldsymbol{10}}$}};
\draw[fill=Couleur10!15, draw=Couleur10, very thick] (2.4,2.2) rectangle (8.0,3.0) node[pos=0.5] {\textcolor{black}{${\color{Couleur10}\boldsymbol{109}}$}};
\draw[fill=Couleur10!15, draw=Couleur10, very thick] (0.0,0.9) rectangle (4.8,1.7) node[pos=0.5] {\textcolor{black}{${\color{Couleur10}\boldsymbol{60}}$}};
\draw[<->,thick, draw=black] (4.8,1.5) -- (8.0,1.5) node[pos=0.5, below] {\textcolor{black}{${\color{Couleur10}\boldsymbol{59}}$}};
\draw[dashed, thick, draw=black] (8.0,0.9000000000000001) -- (8.0,1.7);
\end{tikzpicture}
\item Le problème :

\medskip
Yann participe à un triathlon. Il nage $200$ m, fait ensuite du vélo et ensuite court sur $1\,000$ m. Il a parcouru en tout $4\,200$ m. Quelle distance a-t-il parcourue à vélo ?

\medskip

      Réponse :\\
      - Calcul de la distance parcourue en nage et en course à pied en m : $200+1\,000=1\,200$\\
      - Distance parcourue à vélo en m : $4\,200-1\,200=$ $3\,000$.

\medskip

  se modélise par le schéma suivant :

\medskip

\begin{tikzpicture}[scale=0.8]
\draw[decorate,decoration={brace,amplitude=10pt},xshift=0pt,yshift=0pt,draw=black] (0.0,3) -- node[above=10pt, pos=0.5] {\textcolor{black}{${\color{Couleur10}\boldsymbol{4\,200}}$}} (7.2,3);
\draw[fill=Couleur10!15, draw=Couleur10, very thick] (0.0,2.2) rectangle (2.4,3.0) node[pos=0.5] {\textcolor{black}{${\color{Couleur10}\boldsymbol{200}}$}};
\draw[fill=Couleur10!15, draw=Couleur10, very thick] (2.4,2.2) rectangle (4.8,3.0) node[pos=0.5] {\textcolor{black}{${\color{Couleur10}\boldsymbol{3\,000}}$}};
\draw[fill=Couleur10!15, draw=Couleur10, very thick] (4.8,2.2) rectangle (7.2,3.0) node[pos=0.5] {\textcolor{black}{${\color{Couleur10}\boldsymbol{1\,000}}$}};
\end{tikzpicture}

\columnbreak

\item Le problème :

\medskip
Tania effectue à vélo une étape de montagne du tour de France. Elle grimpe trois cols qui ont les dénivelés suivants : $600$ m pour le premier col, $1\,200$ m pour le deuxième et $950$ m pour le dernier. Quelle dénivelé cumulé a-t-elle grimpé ?

\medskip

      Réponse : Tania a grimpé au total en m : $600+1\,200+950=$ $2\,750$.

\medskip

  se modélise par le schéma suivant :

\medskip

\begin{tikzpicture}[scale=0.8]
\draw[decorate,decoration={brace,amplitude=10pt},xshift=0pt,yshift=0pt,draw=black] (0.0,3) -- node[above=10pt, pos=0.5] {\textcolor{black}{${\color{Couleur10}\boldsymbol{2\,750}}$}} (7.2,3);
\draw[fill=Couleur10!15, draw=Couleur10, very thick] (0.0,2.2) rectangle (2.4,3.0) node[pos=0.5] {\textcolor{black}{${\color{Couleur10}\boldsymbol{600}}$}};
\draw[fill=Couleur10!15, draw=Couleur10, very thick] (2.4,2.2) rectangle (4.8,3.0) node[pos=0.5] {\textcolor{black}{${\color{Couleur10}\boldsymbol{1\,200}}$}};
\draw[fill=Couleur10!15, draw=Couleur10, very thick] (4.8,2.2) rectangle (7.2,3.0) node[pos=0.5] {\textcolor{black}{${\color{Couleur10}\boldsymbol{950}}$}};
\end{tikzpicture}
\item Le problème :

\medskip
Karine a $5$ boîtes de chocolats chacun contenant $30$ chocolats. Quand elle les distribue à ses amis, chacun en a $37$ et il en reste 2. Combien a-t-elle d'amis ?

\medskip

    Réponse :\\
    - Calcul du nombre de chocolats total : $5 \times 30 = 150$.\\
    - Calcul du nombre de chocolats partagés : $150-2=148$.\\
    - Nombre d'amis : $148 \div 37 = 4$.\\
    Karine a $4$ amis.

\medskip

  se modélise par le schéma suivant :

\medskip

\begin{tikzpicture}[scale=0.8]
\draw[decorate,decoration={brace,amplitude=10pt},xshift=0pt,yshift=0pt,draw=black] (0.0,3) -- node[above=10pt, pos=0.5] {\textcolor{black}{${\color{Couleur10}\boldsymbol{5}}$}} (8.0,3);
\draw[fill=Couleur10!15, draw=Couleur10, very thick] (0.0,2.2) rectangle (2.4,3.0);
\node[anchor=center, align=center] at (1.2,2.7) {\textcolor{black}{${\color{Couleur10}\boldsymbol{30}}$}};
\draw[fill=Couleur10!15, draw=Couleur10, very thick] (2.4,2.2) rectangle (8.0,3.0);
\node[anchor=west, align=left] at (2.4,2.7) {\textcolor{black}{\ldots}};
\draw[fill=Couleur10!15, draw=Couleur10, very thick] (0.0,0.9) rectangle (2.4,1.7);
\node[anchor=center, align=center] at (1.2,1.4) {\textcolor{black}{${\color{Couleur10}\boldsymbol{37}}$}};
\draw[fill=Couleur10!15, draw=Couleur10, very thick] (2.4,0.9) rectangle (5.6,1.7);
\node[anchor=west, align=left] at (2.4,1.4) {\textcolor{black}{\ldots}};
\draw[fill=Couleur10!15, draw=Couleur10, very thick] (5.6,0.9) rectangle (8.0,1.7);
\node[anchor=center, align=center] at (6.8,1.4) {\textcolor{black}{$2$}};
\draw[<->,thick, draw=black] (0.0,0.7000000000000002) -- (5.6,0.7000000000000002) node[below, pos=0.5] {\textcolor{black}{${\color{Couleur10}\boldsymbol{4}}$}};
\end{tikzpicture}
\end{enumerate}
\end{multicols}
 

\clearpage

\end{document}

% Local Variables:
% TeX-engine: luatex
% End:


\end{document}